\section{Additional steps in cote 149, scans 105--136}
The following statements record the complete sweep of \S4. The formulas already answered in D1--D10 and E1--E6 are not duplicated. We retain a genuinely unresolved assertion as an explicitly identified obstruction; no conjectural assertion is used in any preceding proof.

\begin{result}{ZD1}{Proved}{A simply connected scheme as a fiber functor}
Let $Y$ be connected and locally noetherian, and let $Z\to Y$ be nonempty and connected with trivial \emph{\'etale} fundamental group. Then
\[
 F_Z:\operatorname{FEt}(Y)\longrightarrow\operatorname{FinSets},\qquad
 E\longmapsto\pi_0(E\times_Y Z)
\]
is a fiber functor. A map of such pointed data gives a genuine continuous homomorphism of its automorphism groups. For a morphism $f:X\to Y$, the relative kernel depends on the choice of connecting path through conjugation by $\pi_1(X)$; that action need not be inner as an action on the kernel itself.
\par\emph{Passages:} cote 149, \S4, manuscript pp.~52--56, scans 105--114, especially ``on devrait pouvoir \`ecrire'' at scan 107.
\end{result}
\begin{proof}
Every finite \emph{\'etale} cover of $Z$ is a disjoint union of copies of $Z$. Choose a geometric point $z$ of $Z$. Each of these copies has exactly one point above $z$, giving a natural bijection $F_Z(E)\simeq E_z$. Hence $F_Z$ is isomorphic to the usual geometric fiber functor, and is exact and faithful. This identification also constructs every claimed comparison map. A different connecting path between two geometric fiber functors is obtained by the action of an element of $\pi_1(X)$, so the comparison of kernels is its restricted conjugation. It is generally an outer automorphism of the kernel: for $X=\mathbb G_{m,\mathbb Q}$ and $Y=\operatorname{Spec}\mathbb Q$, conjugation acts cyclotomically on the abelian kernel $T$, whereas every inner automorphism of $T$ is trivial. Thus the manuscript's qualification on scan 113 is necessary, and ``automorphismes int\'erieurs'' on scan 114 refers to the ambient group.
\end{proof}

\begin{result}{ZD2}{Proved}{Relative kernels, universal covers, and the missing left exactness}
For connected locally noetherian pointed schemes $f:X\to Y$, put $H=\operatorname{im}(\pi_1(X)\to\pi_1(Y))$. The pointed connected pro-finite-\emph{\'etale} cover $Y_H\to Y$ corresponding to $H$ admits a unique pointed lift of $f$, the map $\pi_1(X)\to\pi_1(Y_H)=H$ is surjective, and its kernel is unchanged. Pullback of the universal pro-cover of $Y_H$ has fundamental group that kernel. The identification of this kernel with a geometric fiber group additionally requires injectivity of the fiber homomorphism; an exact homotopy sequence beginning only at that fiber group does not supply it.
\par\emph{Passages:} cote 149, \S4, manuscript pp.~54--56, scans 109--113, formulas (12)--(20).
\end{result}
\begin{proof}
The finite quotients of the coset space $\pi_1(Y)/H$, with its distinguished coset, define $Y_H$. Pullback to $X$ has that coset fixed by $\pi_1(X)$, hence a compatible section at every finite level; this constructs the lift as a morphism to the pro-cover. The homomorphism to $H$ is surjective by its definition, and inclusion $H\hookrightarrow\pi_1(Y)$ is injective, proving equality of the two kernels. Pulling the universal cover back along this surjection gives the connected pro-cover of $X$ corresponding to the kernel. Its fundamental group is that closed subgroup: a finite \emph{\'etale} cover of the inverse limit descends to a finite level because its defining algebra and \emph{\'etaleness} involve finitely many equations, and the usual finite-cover Galois correspondence then gives precisely the finite continuous sets for the subgroup. Pullback over $Y$ instead gives components indexed by $\pi_1(Y)/H$; the chosen pointing selects the component just described. These are the comparison maps in (14)--(18).
For a geometric fiber $X_{\bar y}$ the homomorphism $\pi_1(X_{\bar y})\to\pi_1(X)$ has image the relative kernel whenever the usual homotopy exactness holds. It identifies the two groups precisely when it is also injective. The punctured tautological line bundle $\mathcal O_{\mathbb P^1}(-1)^\times=\mathbb A^2\setminus\{0\}$ over an algebraically closed characteristic-zero field proves that this extra point matters: its base and total space have trivial fundamental group, by the projective-line computation and purity, whereas the fiber has group $T$. The map is $T\to1$. This is also the explicit example in \cite[Remark 9]{DE_BES}. Thus the phrase ``suite exacte d'homotopie universelle'' must include the missing left injectivity for (12), (16), and (19) in the manuscript to have their asserted isomorphism meaning. Over a field the ordinary geometric exact sequence does include this injection, so the application at scan 110 is valid.
\end{proof}

\begin{result}{ZD3}{Proved}{Recovering the model quotient and the function field}
For a regular integral finite-type model $U$ of $L/K$, the kernel of $G_L\to\pi_1(U)$ is the closed normal subgroup generated by the inertia groups of its codimension-one local rings. It is equally generated by the inertia groups of all valuations with a center on $U$. For the nonempty open subschemes of one fixed normal model,
\[
 G_L\simeq\varprojlim_U\pi_1(U),\qquad
 P_{L/K}\simeq\varprojlim_U\pi_1(U_{\overline K})
\]
when the models are geometrically connected; in general the second formula uses the kernels of the maps to $G_K$. For $U=X\setminus Z$ with $X$ proper regular, the kernel of $\pi_1(U)\to\pi_1(X)$ is generated by the boundary-divisor inertia images.
\par\emph{Passages:} cote 149, \S4, manuscript pp.~57--60, scans 116--121, (21)--(25) and the first paragraph of scan 121.
\end{result}
\begin{proof}
A finite Galois quotient of $G_L$ killing the listed codimension-one inertias gives a finite normalization of $U$ unramified in codimension one; D1's purity argument makes it finite \emph{\'etale}. Conversely every such cover kills the inertia of every valuation with a center, because its pullback to a strictly henselian valuation ring is trivial. Taking all finite quotients proves equality of the closed normal subgroups. The same argument applied to the boundary divisors of $X$ proves the last assertion.
Every finite separable extension of $L$ is finite \emph{\'etale} over a sufficiently small nonempty affine open: choose integral equations for field generators, invert their denominators and a nonzero discriminant, and take the resulting finite algebra. Thus the finite quotients in the inverse system exhaust those of $G_L$. Restriction of a connected finite \emph{\'etale} cover of a normal integral scheme to a nonempty open remains connected, because that cover is normal integral. This proves the surjectivity of the transition maps and of each generic-point map. Equality of the finite quotients gives the first isomorphism as an actual isomorphism with the fixed generic fiber functor. Taking kernels of the compatible maps to $G_K$ commutes with the inverse limit and proves the second formula. No section-conjecture input is involved.
\end{proof}

\begin{result}{ZD4}{Refuted}{Model opens, finite unions, and when the Riemann--Zariski space is a scheme}
The points of manuscript codimension one are exactly the divisorial valuation rings, as asserted on scans 121--122. The Riemann--Zariski space is the proper normal curve for transcendence degree one, a point for degree zero, and has no nonempty affine open in its locally ringed-space structure for transcendence degree at least two. The claim on scan 123 that the opens supplied by models are not stable under finite unions is false: these opens are precisely the quasicompact opens, and they are stable under finite unions and finite intersections.
\par\emph{Passages:} cote 149, \S4, manuscript pp.~60--61, scans 121--123.
\end{result}
\begin{proof}
Write $n=\operatorname{trdeg}_K L$. Suppose a nontrivial valuation has residue transcendence degree $n-1$. Choose units $u_1,\ldots,u_{n-1}$ with algebraically independent residues and an element $t$ of positive value. The least-value argument of D4 makes these $n$ elements algebraically independent, so $L$ is finite over $K(u_1,\ldots,u_{n-1},t)$. Its restricted valuation is the $t$-adic valuation with coefficient field $K(u_1,\ldots,u_{n-1})$: independence of the residues ensures that every nonzero rational coefficient has value zero. The fundamental inequality for this finite field extension makes both the value-group index and the residue-field degree finite. The value group is consequently discrete, and the residue field finitely generated of transcendence degree $n-1$. Normalize $K[u_1,\ldots,u_{n-1},t]$ in $L$. The center lies above the prime $(t)$ and has height one; its local normal ring is a DVR, dominated by the given DVR of the same fraction field, hence equal to it. This realizes the valuation as a divisor on a model. Conversely a codimension-one point of a normal finite-type model has exactly this residue transcendence degree by the dimension formula. The valuation of manuscript codimension zero is trivial by D4's inequality. For $n=1$ every nontrivial valuation is the local DVR of its center on the proper normal curve, proving the assertion about the whole locally ringed space by D2.
For the topology, an open of the form $p_X^{-1}(U)$, for a model open $U\subset X$, consists exactly of the valuations with center on $U$. Such opens form a basis and are quasicompact. Indeed, the valuation-ring description embeds the space as a closed subspace of the compact product $\{0,1\}^{L}$ in the constructible topology, by the valuation-ring axioms. For an affine model open $\operatorname{Spec}A$, having a center there says that the finitely many algebra generators of $A/K$ lie in the valuation ring. This condition is clopen in the compact constructible topology, proving the required quasicompactness. A general model open is a finite union of affine ones. A quasicompact open is a finite union of such basic opens. Dominate their finitely many proper models by the closure of the common generic point in their product. On that common model their union, or intersection, is the inverse image of the union, or intersection, of the corresponding model opens. This proves the asserted description and stability, and directly corrects the sentence on scan 123.
Finally suppose a nonempty affine open $W$ existed in the Riemann--Zariski space for $n\ge2$. By what we just proved, $W=p_X^{-1}(U)$ for a nonempty open of a normal proper model. Its ring of sections is
\[
 A=\Gamma(W,\mathcal O)=\bigcap_{V\in W}V=\Gamma(U,\mathcal O_U).
\]
The first equality with the intersection follows from D2's local-ring description: membership of a rational function in a valuation ring is witnessed on a neighborhood of its center on some model. The last equality follows by testing the function at all codimension-one DVRs of the normal scheme $U$, which is the local criterion for regularity on a normal scheme. Pick a smooth closed point $a\in U$. There are distinct valuation rings centered at $a$: blow up $a$, take the valuation of its exceptional divisor, and refine that valuation by a nontrivial valuation of the exceptional divisor's positive-transcendence-degree function field. Both lie in $W$. Both contract their maximal ideals to the same prime $\mathfrak m_a\cap A$, because $A\subset\mathcal O_{U,a}$ and both dominate this local ring. This is impossible in $W=\operatorname{Spec}A$, where a point is exactly its contracted prime. Thus no such affine open exists. This proves the nonscheme claim without resolution of singularities.
\end{proof}

\begin{result}{ZD5}{Refuted}{Nonuniqueness of the valuation normalized by a section}
The uniqueness and the decomposition-group inference in conjecture (a) of scans 126--127 are false for arbitrary valuations, even with nontrivial inertia and residue field exactly $K$. The normalization condition records the coarsening $w$ of D5, and need not recover its residual refinement $v/w$.
\par\emph{Passage:} cote 149, \S4, manuscript pp.~62--63, scans 126--127, ``et ce $V$ est uniquement'' and ``en fait la section provient''.
\end{result}
\begin{proof}
Let $K=\mathbb Q$, $L=\mathbb Q(z,x,y)$, and let $w$ be the $z$-adic valuation. On its residue field $\mathbb Q(x,y)$ use the two valuations $r_0,r_1$ of D3 with value group $\mathbb Q$ and residue field $\mathbb Q$, obtained by substituting $x-a=t$, $y=2+\sum t^{j+1/j!}$ for $a=0,1$. They are distinct because their centers on $(\mathbb P^1)^2$ are $(0,2)$ and $(1,2)$. Compose them with $w$ to obtain $v_0,v_1$. Their value groups are $\mathbb Z\oplus_{\rm lex}\mathbb Q$, their residue fields are $\mathbb Q$, and D5 gives
\[
 I_{v_0}=I_{v_1}=I_w=T\ne1,\qquad N_{G_L}(I_{v_a})=D_w.
\]
There is a section $G_{\mathbb Q}\to D_w$: first construct a section $G_{\mathbb Q}\to G_{\mathbb Q(x,y)}$ by successive decomposition splittings along the rational flag $y=2$, $x=0$, and then use the splitting for $w$ in D4. Its image normalizes both displayed identical inertia groups, proving nonuniqueness with nonzero inertia.
This section is not contained in $D_{v_1}$. Restriction to $\mathbb Q(x)$, followed by the quotient for $C=\mathbb P^1\setminus\{2,3,\infty\}$, is the point section at $0$. Containment in $D_{v_1}$ would make it the point section at $1$, since $r_1(x-1)>0$. Point-section injectivity for this hyperbolic curve, as proved in B and in the prior return's Kummer argument, excludes equality or conjugacy of those sections. Thus actual containment in a decomposition group does not follow from normalization of inertia. Without the nonzero-inertia qualification, the two valuations $r_0,r_1$ themselves already have $I=1$ and are normalized by every section. These examples explain both failures by the explicit homomorphisms, rather than by a change of terminology.
\end{proof}

\begin{result}{ZD6}{Refuted}{The overly broad centralizer margin}
The assertion in the margin of scan 127 that every subgroup of $N_v$ with open image in $N_v/I_v$ has centralizer $I_v$ is false. Cyclotomic triviality is necessary even for inclusion of $I_v$ in that centralizer. For decomposition sections of rational places in a one-variable function field, the corrected cyclotomic-kernel statement is the one proved in D9 and the prior return B7--B8.
\par\emph{Passage:} cote 149, \S4, manuscript p.~63, scan 127, ``ou simplt $\forall$ s-groupe de $N_V$ dont l'image ... soit ouvert''.
\end{result}
\begin{proof}
Take $L=\mathbb Q(t)$ and the $t$-adic valuation. A decomposition section $s:G_{\mathbb Q}\to D_v$ has image projecting onto all of $G_{\mathbb Q}=D_v/I_v$. The conjugation formula of D4 is
\[
 s(g)\tau s(g)^{-1}=\chi(g)\tau\quad(\tau\in T=I_v).
\]
Complex conjugation has $\chi(g)=-1$ and does not centralize any nonzero $\ell$-adic inertia element for an odd prime $\ell$. Therefore $I_v$ is not even contained in the claimed centralizer. For $s(G_{\mathbb Q}^{\circ})$ the action on $I_v$ is trivial, and the reverse inclusion is the independently proved curve result cited in D9. The two assertions cannot be equated without restricting the subgroup in the margin.
\end{proof}

\begin{result}{ZD7}{Refuted}{Boundary inertia images do not determine the valuation or form a model-independent list}
Distinct valuations can have the same inertia image in $\pi_1(U)$, and even the collection of local boundary-inertia images, up to conjugacy and commensurability, changes under a blow-up. Thus the intrinsicness asserted at scan 132 does not follow and is false in that form.
\par\emph{Passages:} cote 149, \S4, manuscript pp.~63--66, scans 128, 131--132.
\end{result}
\begin{proof}
All valuations centered in $U$ have zero inertia image by D7. D3 gives distinct valuations centered outside $U$ with the same zero image. It also gives distinct refinements with exactly equal nonzero inertia subgroups by ZD5. These already disprove injectivity of the assignment $v\mapsto\Delta_v$.
For dependence of the local-image collection, take the product compactification of $U=(\mathbb P^1\setminus\{0,1,\infty\})^2$ and blow up the crossing $(0,0)$. At a general point of the exceptional divisor use $x=u$, $y=uv$ with $v$ a nonzero finite unit. Its meridian maps into the original cusp product $T_x\times T_y$ by $a\mapsto(a,a)$. On the original compactification the nonzero local images are a cusp group in one curve factor or a full product of two cusp groups. A diagonal group has nontrivial projections into both factors, so cannot be conjugate or commensurable with a one-factor group. It has infinite index in a two-factor product: already on each $\ell$-adic factor the quotient has rank one. Thus the blow-up creates an image not commensurable with any original local image. The homomorphism is the Kummer map $v(x)=v(y)=1$ of E3, which verifies every comparison used here.
\end{proof}

\begin{result}{ZD8}{Proved}{What the boundary incidence complex does describe}
For a smooth complex variety, the incidence complex of a simple normal crossings boundary has the same connected components as the boundary. It must have a cell for every connected component of each intersection, not merely one simplex for each nonempty set of labels. It does not in general give the homotopy type of a punctured neighborhood, even when all its strata are points. The pro-system of complements of compact subsets of $U(\mathbb C)$ is independent of its smooth proper compactification; punctured boundary neighborhoods are a cofinal system in it.
\par\emph{Passage:} cote 149, \S4, manuscript p.~66, scan 133, ``quelle est sa signification topologique?'' and the withdrawn contractibility proposal.
\end{result}
\begin{proof}
A connected component of the boundary is precisely a maximal set of its connected smooth components joined by a chain of nonempty intersections. This is exactly connectivity of the incidence graph and hence of the incidence complex. To see why intersection components must be retained, in $\mathbb P^1\times\mathbb P^1$ take the graphs $y=x$ and $y=1/x$. They meet transversely at $(1,1)$ and $(-1,-1)$. Their actual incidence complex has two vertices joined by two edges, whereas a single simplex per label set gives only one edge and loses that circle.
For $U=\mathbb P^1(\mathbb C)\setminus\{0,1,\infty\}$ all boundary strata are points. The incidence complex is three discrete points. A small punctured neighborhood of the boundary is three punctured discs, each with fundamental group $\mathbb Z$; it cannot have that homotopy type. The torus factors and their Chern gluing in E4 are precisely the additional data in this example and its higher-dimensional counterparts.
Finally, for any compact subset $C\subset U(\mathbb C)$, the open set $X(\mathbb C)\setminus C$ contains the boundary. Conversely the complement of a boundary neighborhood is a closed subset of the compact space $X(\mathbb C)$ disjoint from the boundary, hence a compact subset of $U(\mathbb C)$. These inclusions identify the two cofinal systems directly and prove the stated intrinsic end pro-system. The connected components of sufficiently small punctured tubular neighborhoods agree with those of the boundary, since the normal punctured-disc fibers are connected and tubes over intersecting branches meet. We assert this exact comparison, not an identification with the bare incidence complex.
\end{proof}

\begin{result}{ZD9}{Proved}{Conjugacy along strata and the arithmetic stabilizer}
On a geometrically connected open stratum with labelled geometric boundary branches, its homomorphisms $T^J\to\pi_1(U_{\overline K})$ are conjugate by geometric paths. Arithmetic transport acts on the labels by permutation and on each $T$ by the cyclotomic character. An individual geometric branch carries canonical equivariance under its Galois stabilizer; the full Galois group preserves the orbit of branches, and invariance of an individual image can fail. Thus scan 136's full-group assertion needs geometric connectedness or this stabilizer correction.
\par\emph{Passages:} cote 149, \S4, manuscript p.~67, scans 135--136, (38)--(40).
\end{result}
\begin{proof}
The local Kummer constructions of E1 glue along the normal-bundle torsor of E4. A path in the connected geometric stratum identifies its two torus fibers. A lift of this path to the total space conjugates their fundamental-group homomorphisms; a second lift changes the conjugator by an element of the torus fiber, which is abelian and acts trivially on its own group. Composition with the morphism to $U$ gives the claimed geometric conjugacy, whether or not the meridian map is injective. On geometric fibers the branch labels are fixed. An arithmetic path permutes the geometric branches and sends a root of unity $\zeta_n$ to $\zeta_n^{\chi(g)}$, so its action is exactly the stated permutation and cyclotomic action. This proves the action on actual finite Kummer covers, and passing to their inverse limit proves the statement for $T^J$.
For a concrete failure without the correction, take $U=\mathbb P^1_{\mathbb Q}\setminus\{\sqrt2,-\sqrt2,\infty\}$. The degree-two closed boundary point is irreducible and connected over $\mathbb Q$, but has two geometric branches. An element of $G_{\mathbb Q}$ interchanging them interchanges their cusp subgroups in the geometric free profinite group of rank two. These subgroups are not conjugate, by the distinct-cusp theorem of A. Thus neither individual branch homomorphism is fixed up to geometric conjugacy by all of $G_{\mathbb Q}$. Its stabilizer is $G_{\mathbb Q(\sqrt2)}$, and the two-element orbit has the full equivariance claimed above.
\end{proof}

\begin{result}{ZD10}{Open}{The remaining birational lifting assertion and a smallest obstruction}
The following exact assertion from conjecture (b) is not proved in this return: for $K=\mathbb Q$ and $U=\mathbb P^1\setminus\{0,1,\infty\}$, every continuous section of $\pi_1(U)\to G_K$ lifts through the continuous quotient $G_{K(t)}\to\pi_1(U)$ to a continuous section of $G_{K(t)}\to G_K$. Likewise, after the false uniqueness clause in conjecture (a) is removed, the assertion that every section of $G_{K(t)}\to G_K$ is contained in a rational-place decomposition group remains unproved here. These unproved assertions are not premises of D or E.
\par\emph{Passages:} cote 149, \S4, manuscript pp.~62--63, scans 126--127, conjectures (a) and (b).
\end{result}
\noindent\emph{Attack and exact point of failure.}
All point sections and all cuspidal sections in this example do lift: at an interior rational point or a rational cusp, the decomposition sequence has the continuous splitting of D4, and its image under $G_{K(t)}\to\pi_1(U)$ is precisely the point or tangential section. For an arbitrary section $s$, form the actual profinite pullback
\[
 E_s=G_{K(t)}\times_{\pi_1(U),s}G_K,
 \qquad 1\longrightarrow\ker(G_{K(t)}\to\pi_1(U))\longrightarrow E_s\longrightarrow G_K\longrightarrow1.
\]
The required lift is exactly a continuous splitting of this extension. ZD3 describes its kernel by the interior-point inertia normal closure, but that description neither produces a complement nor kills the finite embedding obstructions of $E_s$. The normalizer attack identifies decomposition containment for a valuation already supplied, by D5 in rank one; it gives no existence of a normalized inertia subgroup for an arbitrary section. Passing to the abelianized kernel retains only the abelian lifting obstructions and supplies no implication back to a splitting of $E_s$. Thus these are exact remaining obstructions, not a proof by linearization. The source search checked the opened Borne--Emsalem--Stix theorem: its Theorem 27 lifts sections over $\mathbb Q$ for positive-genus proper curves and divisors that are unions of torsion packets to a specified cuspidally central quotient. That quotient is explicitly smaller than $G_{K(t)}$, its base curve has positive genus, and it supplies no splitting of the displayed $E_s$ for the thrice-punctured line. The attached prior return also records these cases as unclosed. No claimed universal lift or arbitrary-section classification follows from those results.

\begin{result}{ZD11}{Known}{The universal lifting claim has the strength of the section conjecture}
For a fixed smooth geometrically connected hyperbolic curve $C/k$, with $k$ finitely generated over $\mathbb Q$, the following two universal assertions are equivalent: every section of $C_F/F$ is geometric or cuspidal for every finitely generated extension $F/k$; and every such section lifts through $G_{F(C)}\to\pi_1(C_F)$. Thus even the curve case of the unrestricted conjecture (b) in scan 127 contains the full section-conjecture assertion.
\par\emph{Passage:} cote 149, \S4, manuscript p.~63, scan 127, conjecture (b).
\end{result}
\begin{proof}
This is Bresciani, Theorem B, with its hypotheses exactly as stated above \cite{DE_Bresciani}. The comparison underlying its use here is explicit. For an indeterminate $t$, geometric invariance of the curve fundamental group identifies
\[
 \pi_1(C_{F(t)})\simeq\pi_1(C_F)\times_{G_F}G_{F(t)}.
\]
The base change of a section $s$ is the continuous map $g\mapsto(s(\operatorname{res}(g)),g)$ into this fiber product. Bresciani's Theorem A states, for a smooth geometrically connected curve over a field finitely generated over $\mathbb Q$, that $s$ is geometric or cuspidal exactly when this base-changed section is birationally liftable. Universal lifting applies in particular to the finitely generated extension $F(t)/k$, so that theorem gives the first universal assertion. Conversely a geometric section lifts using the rational-place decomposition sequence of D4; a cuspidal section already factors through that sequence, and its identical local splitting defines the birational lift. Quotienting by the geometric conjugacy ambiguity and lifting a conjugator through the surjection $G_{F(C)}\to\pi_1(C_F)$ makes this an assertion about actual continuous lifts as well as section classes. These maps prove the claimed equivalence and locate the unresolved step of ZD10 exactly.
\end{proof}
